% Generated draft: requires mathematical and source review.
\documentclass[11pt]{article}
\usepackage[margin=25mm]{geometry}
\usepackage{fontspec}
\setmainfont{Latin Modern Roman}
\usepackage{amsmath,amssymb,mathtools,unicode-math}
\setmathfont{Latin Modern Math}
\usepackage{enumitem,needspace,xurl,hyperref}
\setlength{\parindent}{0pt}
\setlength{\parskip}{5pt}
\setlength{\emergencystretch}{3em}
\newcommand{\N}{\mathbb{N}}
\newcommand{\Z}{\mathbb{Z}}
\newcommand{\Q}{\mathbb{Q}}
\newcommand{\R}{\mathbb{R}}
\newcommand{\C}{\mathbb{C}}
\newcommand{\F}{\mathbb{F}}
\newcommand{\A}{\mathbb{A}}
\newcommand{\PP}{\mathbb{P}}
\newcommand{\E}{\mathbb{E}}
\newcommand{\eps}{\varepsilon}
\newcommand{\dd}{\,\mathrm{d}}
\begin{document}
\section*{UnsolvedMath ID: 20002046}
\textbf{Source problem code:} catalogue entry
\subsection{Definitions and mathematical statement}
Let $M$ be a closed smooth manifold of even dimension $n$, with a Riemannian conformal class. Treat critical $Q$-curvature as the density $q_g=Q_n(g)\,\dd V_g$. For $\widehat g=e^{2\omega}g$, use the critical transformation law
\[
 q_{\widehat g}=q_g+(P_g\omega)\,\dd V_g,
 \qquad (P_{\widehat g}f)\,\dd V_{\widehat g}
                 =(P_gf)\,\dd V_g.
\]
Here $P$ is the critical GJMS operator, with $P_g1=0$. Define the unmodified functional suggested by Tom Branson,
\[
 H[\widehat g,g]=\int_M\omega(q_{\widehat g}+q_g),
 \qquad \omega=\log\Omega.
\]
The source's displayed $\widehat Q+Q$ must thus use densities or the corresponding separate volume factors; adding two unweighted scalars with no integration convention would not specify this functional.

\textbf{Problem 24(b).} Do Weyl structures help with the preceding, unspecified issue? The extract proposes $H$ and asserts that its three-metric cocycle property is equivalent to self-adjointness of $P$. The antecedent ``in this regard'' is missing, so no broader objective is inferred.

\textbf{Supplied exact cocycle test.} For $g_0=g$, $g_1=e^{2\omega}g$, $g_2=e^{2(\omega+\phi)}g$, the aid gives
\[
 H[g_2,g_1]+H[g_1,g_0]-H[g_2,g_0]
   =\int_M(\phi P_g\omega-\omega P_g\phi)\,\dd V_g.
\]
Consequently the cocycle identity for all conformal triples is equivalent to formal self-adjointness, namely symmetry of this pairing for all smooth $\omega,\phi$. If tests are restricted to a subspace, they detect only the restricted skew pairing. In particular every one-dimensional conformal-ray test, $\omega=af$ and $\phi=bf$, has zero defect for every linear $P$ and cannot establish full self-adjointness.

\textbf{Supplied Weyl-scale descent criterion.} A Weyl connection is torsion free and preserves the conformal class; in a scale $g$ write $Dg=-2\theta\otimes g$. The connection induced by a metric is its Levi-Civita connection, an exact Weyl connection. Two conformally related metrics induce the same connection precisely when their logarithmic scale ratio is constant on each connected component. For the actual self-adjoint critical GJMS operator, let
\[
 \mathcal Q_j=\int_{M_j}q_g
\]
for each connected component $M_j$; these numbers are conformally invariant. Changing the representatives of a pair by $g_0\mapsto e^{2c_j}g_0$ and $g_1\mapsto e^{2d_j}g_1$ changes the functional by
\[
 2\sum_j(d_j-c_j)\mathcal Q_j.
\]
The aid therefore asserts that this unmodified $H$ descends to a function of pairs of metric-induced Weyl connections exactly when $\mathcal Q_j=0$ for every component. On a connected manifold this is simply zero total $Q$; on a disconnected manifold cancellation of the global sum is insufficient. This is a criterion for exact metric-induced Weyl structures, not a construction on all nonexact Weyl connections, nor a claim that no modified functional could remove the obstruction. Boundary versions require boundary conditions and additional terms and are not included.
\subsection{Short English statement}
Branson's proposed functional is a cocycle exactly when the critical operator is formally self-adjoint, provided all conformal triples are tested. The supplied aid also shows that it depends only on metric-induced Weyl connections precisely when total $Q$ vanishes on each connected component. The missing broader aim and nonexact Weyl case remain unspecified.
\subsection{Sources}
\begin{itemize}
\item Local UnsolvedMath registry, permanent ID 20002046, source code catalogue entry; canonical \texttt{aim-geometry-notes.json}, notes[383].
\item AIM \emph{Conformal structure in geometry, analysis, and physics}, Problem 24(b), functional suggested by Tom Branson: \url{https://aimath.org/WWN/confstruct/confstruct.pdf}.
\item No separate bibliographic citation for the cocycle calculation or Weyl-scale descent criterion is supplied; the latter is the research aid's partial result.
\item Registry research attempt 1, dated 2026-08-14: partial result, rigor ``proved'', novelty confidence low. The explicitly machine-generated, non-peer-reviewed aid requires independent verification before citation.
\end{itemize}
\end{document}
